% !TEX root = ../main.tex

\chapter{全文总结与未来展望}

\section{全文工作总结}

本文以 \textit{visualnav-transformer} 项目为基础，以 NoMaD 为高层核心算法，以云深处 Lite3 四足机器人为主要部署对象，围绕“基于动作扩散策略的足式机器人视觉目标导航与探索”展开了从算法到系统再到平台的完整研究。与只停留在离线指标上的工作不同，本文始终把研究目标限定为“算法、系统与平台能否形成可运行闭环”。

围绕这一目标，本文主要完成了四方面工作。第一，系统梳理了 NoMaD 的代码结构、训练目标与推理流程，明确了视觉编码器、距离预测头和扩散策略头之间的共享条件空间关系。第二，围绕部署需求对高层推理进行了优化，重点研究了 DDIM 推理加速、CFG 引导增强、TTS 推理时搜索与视觉编码器升级，建立了统一的离线统计实验链路。第三，面向 Lite3 构建了由统一推理模块、状态机、导航主机、中层控制映射与平台接口组成的分层闭环系统，并在 MuJoCo 中完成了系统级验证。第四，形成了面向 Jetson Orin 与 Lite3 真机平台的完整部署流程，并为真实实验的图像记录、可视化和任务组织留出了规范化接口。

从阶段性结论看，DDIM-2 是当前最适合进入闭环验证的推理配置；CFG 更适合作为可调部署参数，而不是默认固定选项；TTS 已经能够被整合进统一推理模块，并作为推理外层搜索机制继续开展统计实验；视觉编码器升级工作已经完成第一阶段统一筛选框架，但“最终部署最优编码器”的结论仍需第二阶段训练和闭环验证进一步确认。

\section{主要贡献}

本文的主要贡献可归纳为以下三点：
\begin{enumerate}
    \item 建立了面向部署的 NoMaD 推理优化链路，将 DDIM、CFG、TTS 和视觉编码器替换纳入同一推理接口进行分析与控制；
    \item 构建了适用于 Lite3 的分层闭环系统，用状态机与导航主机把高层 NoMaD 策略组织成可恢复、可切换、可扩展的任务链路；
    \item 给出了 Orin 上的 Lite3 真机部署方案，并将图像记录、可视化与安全控制纳入统一工作流，提升了系统复现性。
\end{enumerate}

\section{研究不足}

本文仍存在三方面不足。第一，视觉编码器实验目前主要完成了第一阶段统一筛选，尚未完成面向最终部署结构的 fully-tuned 第二阶段训练，因此对 DINOv2-Small 等候选的结论仍应保持谨慎。第二，真机实验的正式统计表和大规模重复实验尚未全部补齐，因此第五章仍保留了一部分结果占位。第三，当前 topomap 仍主要依赖预先采集的视觉节点，尚未形成更强的在线更新与自维护能力。

\section{未来展望}

后续工作可从以下方向继续推进：
\begin{enumerate}
    \item 完成视觉编码器第二阶段训练与公平闭环验证，给出最终部署推荐；
    \item 补齐真机统计实验，建立离线指标与真机任务成功率之间的更明确对应关系；
    \item 研究可在线更新的 topomap 与视觉记忆机制，降低人工采集成本；
    \item 将当前统一推理模块与状态机扩展到 Tron1 等更多平台，验证分层结构的可迁移性；
    \item 引入语言目标与多模态条件，使系统从 image-goal 进一步扩展到更灵活的任务表达。
\end{enumerate}

总体而言，本文已经把原本主要面向轮式平台验证的 NoMaD 高层视觉导航策略，扩展为一套面向 Lite3 四足平台的可分析、可仿真、可部署的研究链路。后续若能继续补齐第二阶段训练与真机统计证据，该路线有望进一步发展为更稳定的足式视觉目标导航方案。
